\documentclass[11pt]{article}

\usepackage[margin=1in]{geometry}
\usepackage{amsmath,amssymb,amsthm}
\usepackage{booktabs}
\usepackage{multirow}
\usepackage{graphicx}
\usepackage{hyperref}
\usepackage{natbib}
\usepackage{xcolor}
\usepackage{enumitem}
\usepackage{caption}
\usepackage{array}

\hypersetup{colorlinks=true, linkcolor=blue!50!black, citecolor=blue!50!black, urlcolor=blue!50!black}

\newcommand{\risk}{\mathrm{risk}}
\newcommand{\haz}{\mathrm{haz}}
\newcommand{\KL}{\mathrm{KL}}
\newcommand{\E}{\mathbb{E}}

\title{Recurrence-Risk Decoding for Repetition-Loop Control:\\
\large A Minimum-Distortion, Reliability-Analysis Treatment}

\author{}
\date{}

\begin{document}
\maketitle

\begin{abstract}
Autoregressive language models occasionally enter repetition loops: a short span is generated and then reproduced until the sample degenerates into a fixed cycle. This paper studies \emph{recurrence-risk decoding}, a family of decoders that penalize candidate tokens in proportion to the risk that they extend a repeated $n$-gram, and draws a precise distinction between two kinds of control. \emph{Local surrogate control} bounds a one-step proxy for loop risk, re-solved at every token; \emph{path-level event control} reasons jointly over the whole sequence about the probability that the loop event occurs. Every decoder in the family, including the one with an exact proof, is local surrogate control.

The theory is a minimum-distortion (KL-projection) derivation: the exponentially tilted distribution $q(v)\propto p(v)\exp(-\lambda\,\risk(v))$ solves a risk-bounded one-step projection exactly only when $\lambda$ is calibrated per step by solving the dual problem (\textsc{dual} mode), not when $\lambda$ is a fixed constant (\textsc{fixed}) or set heuristically (\textsc{adaptive}). Against an exact path-level oracle, constructed on a finite-state problem through the exponential-tilting solution of a KL-constrained hitting-time problem, even a local surrogate that is given the true multi-step hazard as its risk signal spends 1.6--3.3$\times$ the oracle's KL for the same loop-suppression level (mean 2.2$\times$), and a plain local surrogate spends 2.3--5.1$\times$ (mean 3.1$\times$). The oracle's one-step marginal has a closed form that reweights by hazard affinely rather than exponentially, which explains why an exact risk signal does not close the gap.

Empirically, on a character-level model with a validated persistent-cycle event, \textsc{dual} eliminates the event (0/150 loops) at exactly 100.0\% per-step feasibility, at 11.2--12.7 bits of KL; \textsc{fixed} reaches 2.7\% loops at 0.037 bits and \textsc{adaptive} 0.0\% at 0.096 bits, without a formal guarantee. Against reimplementations of FSD and the LZ Penalty \citep{ginart2025lzpenalty}, the cheapest external baseline that reaches zero loops on this model is LZ Penalty (0.186 bits). The per-step feasibility guarantee holds unchanged on every further model tested (a second checkpoint, a non-fiction-trained model, GPT-2, OLMo-2-0425-1B, and Mistral-7B-v0.3), whereas the effectiveness of every heuristic varies with the model: \textsc{fixed} ranges from 2.7\% to 51.3\% loops across corpora, and \textsc{adaptive} cuts OLMo's loop rate from 12.0\% to 2.0\% but is indistinguishable from noise on Mistral. A blinded human evaluation with 9 participants finds no reliable readability difference among the zero-loop decoders. Transferring the hazard-aware mechanism to a real model under a rollout-based estimator yields no detectable signal.
\end{abstract}

\tableofcontents

%==============================================================================
\section{Introduction}
%==============================================================================

A recurring failure mode of autoregressive text generation is the repetition loop: the model emits a short span and then reproduces it until the sample is dominated by a fixed cycle. It is most visible under low-temperature or deterministic decoding, but the underlying mechanism, the model assigning high probability to continuing a pattern it has already produced, is present at every scale. It is a decoding-time phenomenon: the same trained weights can yield fluent text or a degenerate loop depending only on how the next token is chosen.

The central contribution of this paper is a distinction, made exact and then measured, between two things a decoder can do about loops. Local surrogate control re-solves a one-step proxy for loop risk at every token; every decoder run here does this, including the one with an exact proof. Path-level event control reasons over the whole generated sequence about the actual probability that the loop event ever occurs; this paper does not solve it for a real model. Section~\ref{sec:theory} proves an exact result about the first kind of control. Section~\ref{ssec:oracle} constructs the second kind exactly on a problem small enough to solve, and measures how much the first gives up: 1.6--5.1$\times$ the KL for the same loop-suppression level, even when the surrogate is given the true multi-step hazard. Which particular surrogate mechanism a practitioner should use is a narrower, downstream question that the remaining experiments address.

Every number reported is one of three kinds: an exactly computed property of a well-specified optimization problem; an empirical measurement with a stated sample size and confidence interval, traceable to a saved script; or an explicitly flagged proxy. Implementation defects found and corrected during this work, and their effect on reported numbers, are recorded in Appendix~\ref{app:defects}.

\subsection{Contributions}
\begin{enumerate}[leftmargin=*]
  \item \textbf{The distinction between local surrogate and path-level event control}, with an exact measurement of the gap (Sections~\ref{sec:theory}, \ref{ssec:oracle}). An exact sequence-level oracle shows that local surrogate control spends 1.6--5.1$\times$ the oracle's KL at the same loop-suppression level, even with the exact multi-step hazard as its risk signal, and a closed-form identity for the oracle's one-step marginal (verified to $<10^{-14}$) explains why.
  \item \textbf{A minimum-distortion treatment of recurrence-risk decoding as local surrogate control}: a KL projection with \emph{dual}, \emph{fixed}, and \emph{adaptive} modes, an exact dual-calibration procedure, and its edge cases (Section~\ref{sec:theory}). The guarantee is a statement about the surrogate problem, not the sequence-level event.
  \item \textbf{A validated, deterministic loop event} with a pre-declared calibration rule and a per-corpus false-positive audit (Section~\ref{sec:event}).
  \item \textbf{A controlled comparison} against repetition penalty, hard no-repeat blocking, a suffix-match heuristic, and reimplementations of FSD and the LZ Penalty, including a common-distortion comparison in which every decoder is scored against one shared reference distribution and risk definition (Section~\ref{sec:results}), and a blinded human evaluation (Section~\ref{sec:human-eval}).
  \item \textbf{A synthetic finite-state validation} of hazard-aware local control against the exact oracle (Section~\ref{sec:synthetic}), and a scope-limited null result for transferring it to a real model (Section~\ref{sec:hazard-transfer}).
  \item \textbf{Generalization experiments} on a second character-level model and corpus, GPT-2, and three modern pretrained families (Section~\ref{sec:generalization}), showing that the per-step feasibility guarantee transfers while heuristic effectiveness is model-dependent.
\end{enumerate}

\subsection{Scope}
I do not claim that recurrence-risk decoding eliminates repetition loops in general; that dual calibration provides a sequence-level guarantee against looping; that hard no-repeat blocking is a fair comparator in the primary ranking; that the hazard-aware mechanism transfers from the synthetic problem to a real model; or that any preference ranking among the zero-loop decoders exists or generalizes beyond the single temperature, primary model, and informally recruited participant pool used in the human evaluation.

%==============================================================================
\section{Related Work}
\label{sec:related}
%==============================================================================

\textbf{Decoding strategies.} Temperature sampling rescales logits by $1/T$ before the softmax. Nucleus sampling \citep{holtzman2020curious} restricts sampling to the smallest set of tokens whose cumulative probability exceeds $p$. Repetition penalty \citep{keskar2019ctrl} divides the logit of every previously generated token by a constant, a global discount that does not condition on whether a repeated $n$-gram is imminent; it is the penalty-decoding baseline used throughout. Typical sampling \citep{meister2023typical} keeps tokens whose surprisal is close to the distribution's entropy, and Mirostat \citep{basu2021mirostat} adjusts a truncation threshold online to hold surprisal near a target. No-repeat-$n$-gram blocking \citep{paulus2018deep} zeroes the probability of any token that would complete an already-seen $n$-gram.

\textbf{History-aware and contrastive decoding.} \textsc{SuffixMatchDecoder} is a longest-suffix-match heuristic inspired by LZ77 compression; it is a homemade baseline and is not attributed to any prior publication. Frustratingly Simple Decoding (FSD) \citep{yang2023frustratingly} contrasts the model's prediction against a training-free $n$-gram predictor built from the generated text (the anti-LM) and suppresses candidates the anti-LM is confident about. The FSD-style baseline here is implemented directly from the paper's equations: eq.~2 combines the two in probability space, $p_\theta(v)-\alpha\,p_\omega(v)$, over a top-$k{=}6$ candidate shortlist, with the smoothed multi-order anti-LM of eq.~1/3, $N{=}3$, $\beta{=}0.9$, $\alpha{=}3$ (all the paper's stated defaults). No official code release was available, so this is a faithful implementation of the published formulas, not a certified reproduction of the authors' code, and it is labeled ``reimpl.'' in every table. The LZ Penalty \citep{ginart2025lzpenalty} penalizes each candidate token by the change in codelength it would induce under an LZ-style compressor applied to a recent buffer against an earlier lookback window. The implementation is ported directly from the authors' released reference code (github.com/tginart/sglang) and reproduces it exactly (maximum difference 0.0) on five constructed test cases. More broadly, adaptive contrastive decoding contrasts a strong model's output against a weaker reference; the FSD-style baseline is a repetition-focused instance of that idea, not a general implementation.

\textbf{Survival analysis for degeneration.} A sample that never loops within the observation window is a right-censored observation, not an outlier. Loop onset is therefore treated as a survival time, estimated with Kaplan--Meier \citep{kaplan1958nonparametric} and summarized by Restricted Mean Survival Time (RMST).

\textbf{Constrained decoding as projection.} The dual calibration in Section~\ref{sec:theory} applies Lagrangian duality for convex programs \citep{boyd2004convex} to a KL objective with a linear expectation constraint, which has a closed-form exponential-tilting solution. The character-level models are decoder-only Transformers \citep{vaswani2017attention}.

%==============================================================================
\section{Formal Problem Definition}
\label{sec:formal}
%==============================================================================

Let $p(\cdot\mid x_{<t})$ denote the model's distribution over the vocabulary $\mathcal{V}$ at step $t$. A \emph{decoder} produces a sampling distribution $q_t$ at each step. Greedy decoding sets $q_t=\delta_{\arg\max p}$; temperature sampling sets $q_t(v)\propto p(v)^{1/T}$. Every KL divergence reported is measured against the temperature-scaled model distribution, $p=\mathrm{softmax}(\text{logits}/T)$ at the temperature of that configuration, not the raw $T{=}1$ distribution. At $T{=}0.10$, $p$ is already sharpened and low-entropy, so a KL against it measures movement within a distribution that leaves little room to move cheaply.

\textbf{The loop event.} Fix a detector $\mathcal{L}:\mathcal{V}^{*}\to\{0,1\}$. The \emph{loop onset} of a generation $x_{1:N}$ is $\tau=\min\{t\le N:\mathcal{L}(x_{1:t})=1\}$, or $\tau=\infty$ (right-censored) if no such $t$ exists. Section~\ref{sec:event} develops and validates $\mathcal{L}$.

\textbf{Recurrence risk.} For a candidate $v$ and $n$-gram orders $[n_{\min},n_{\max}]$,
\[
  \risk(v)=\frac{1}{N}\sum_{n=n_{\min}}^{n_{\max}}\mathbb{1}[(\text{last } n{-}1 \text{ tokens},v)\text{ has occurred earlier}],\quad N=n_{\max}-n_{\min}+1.
\]

\textbf{Two levels of control.} The object of interest is $P_Q(\tau\le H)$ under the sampling policy $Q=(q_1,q_2,\ldots)$: the probability that the loop event occurs within a horizon $H$. A policy that minimizes $\KL(Q\Vert P)$ subject directly to a bound on $P_Q(\tau\le H)$, reasoning over the whole path, performs path-level event control. The theorem in Section~\ref{sec:theory} does not solve it, and it is computed exactly here only on a finite-state problem (Section~\ref{ssec:oracle}). Every decoder run in this paper, \textsc{rr-fixed}, \textsc{rr-adaptive}, and \textsc{rr-dual} alike, performs local surrogate control: at each step, independently of every other step, it bounds or penalizes a one-step proxy ($\risk(v)$ for the recurrence-risk family; the exact multi-step hazard $\haz(v)$ for the hazard-aware extension of Section~\ref{sec:synthetic}) and re-solves from scratch at the next step, without accounting for the distortion its decision forces later. The dual mechanism's exact per-step feasibility guarantee is a proof about local surrogate control and is not in dispute. The central quantitative finding (Section~\ref{ssec:oracle}) is that being the exact solution of the surrogate problem is not the same as being close to a solution of the path-level problem the surrogate stands in for.

%==============================================================================
\section{Loop-Event Validation}
\label{sec:event}
%==============================================================================

\subsection{Why the naive event fails}
\label{ssec:naive-fails}
The candidate event ``the first repeated $n$-gram of length $n$'' fires on 54.0\% of held-out Doyle windows and 72.0\% of held-out Dickens windows at $n{=}10$. An event that fires on the majority of ordinary prose cannot distinguish degeneration from fluency and is not used.

\subsection{Calibration}
The persistent-cycle event is the earliest position $e$ at which a span of length $P\in[2,P_{\max}]$ repeats exactly $R$ times consecutively. $R$ and $P_{\max}$ are chosen by a pre-declared rule (the smallest $R$ whose worst-case upper 95\% CI, across window lengths 300/500/1000/2000, falls at or below a 1\% tolerance) on a calibration set (Carroll, Twain, Swift) never referenced afterward. No cell meets the 1\% tolerance at every window length; $R{=}6$ reproduces $R{=}5$'s counts exactly, indicating that the residual is stable and not a threshold-boundary artifact. Frozen: $R{=}5$, $P_{\max}{=}60$, with worst-case upper CI 2.80\%, which does not meet the a priori target.

Period 1 (a single repeated character) uses a separate, decoupled threshold of 8. Using $R{=}5$ for period 1 caused plain-text typesetting conventions (five-character dash dividers) to fire; the threshold of 8 sits above ordinary divider and ellipsis lengths and below observed degenerate runs. Every other period $P\in[2,60]$ uses $R{=}5$.

\subsection{Validation battery}
\begin{table}[h]
\centering
\caption{Validation battery for the frozen event.}
\label{tab:validation}
\begin{tabular}{@{}lrrl@{}}
\toprule
Corpus & fires/$n$ & rate & 95\% CI \\
\midrule
Melville, \emph{Moby-Dick}      & 0/300  & 0.0\% & [0.0\%, 1.2\%] \\
Stoker, \emph{Dracula}          & 26/300 & 8.7\% & [5.7\%, 12.4\%] \\
Whitman, \emph{Leaves of Grass} & 0/300  & 0.0\% & [0.0\%, 1.2\%] \\
Paine, \emph{Common Sense}      & 1/249  & 0.4\% & [0.0\%, 2.2\%] \\
Homer, \emph{Iliad} (trans.)    & 0/300  & 0.0\% & [0.0\%, 1.2\%] \\
Doyle, \emph{Sherlock Holmes}   & 0/300  & 0.0\% & [0.0\%, 1.2\%] \\
Dickens, \emph{A Tale of Two Cities} & 0/300 & 0.0\% & [0.0\%, 1.2\%] \\
\midrule
\textbf{Pooled}                 & 27/2049 & \textbf{1.3\%} & \textbf{[0.9\%, 1.9\%]} \\
\bottomrule
\end{tabular}
\end{table}
The residual is concentrated almost entirely in \emph{Dracula} (epistolary formulas); poetry (Whitman) shows zero false positives. Robustness to window length (300/500/1000/2000 characters) gives 0.0\%/0.5\%/1.4\%/2.4\%.

\subsection{Synthetic sensitivity}
Constructed text with a known period $P\in\{1,2,3,5,8,13,20,30,45,60\}$ repeated $R_{\text{true}}\in\{2,3,4,5,8,10\}$ times is detected correctly in 98--100\% of cases at every tested repeat count.

%==============================================================================
\section{Experimental Protocol}
\label{sec:protocol}
%==============================================================================

\subsection{Model and data}
The primary model is an 848K-parameter, 4-layer, 128-dimensional character-level Transformer with a 256-character context, trained on Austen and Shelley, validated on Doyle, and tested on Dickens. All primary-model results use the \textsc{cosine}-schedule checkpoint unless stated otherwise. A second, independently trained checkpoint of the same architecture and data (\textsc{baseline}, constant learning rate) provides a replication: greedy decoding loops at 60.0\% (vs.\ 93.3\% for \textsc{cosine}) and $T{=}0.15$ still shows a 4.0\% loop rate where \textsc{cosine} shows 0.0\%, so the temperature threshold does not carry over even across retrains of the same architecture and data. \textsc{rr-dual} ($\varepsilon{=}0.05$) reaches 0/150 loops on \textsc{baseline} at 100.0\% feasibility and 11.343 bits of KL, close to \textsc{cosine}'s 11.186 bits (Table~\ref{tab:main}).

\subsection{Where the pathology occurs}
A sweep of 14 temperatures $\times$ 5 prompts $\times$ 10 seeds (1000-character generations, exact Clopper--Pearson intervals) shows the loop rate falling smoothly from 96.0\% at $T{=}0.05$ to 0.0\% by $T{=}0.30$, crossing below 5\% near $T{\approx}0.2$. Every primary-model decoding comparison runs at $T{=}0.10$. The sweep was not repeated for the second model (Section~\ref{sec:generalization}), whose safe temperature range does not simply transfer.

\subsection{Prompt protocol}
A \textsc{Dev} set of 5 prompts is the only material inspected while choosing hyperparameters. A disjoint \textsc{Test} set of 15 prompts, written afterward and not inspected during tuning, is the source of every number in Sections~\ref{sec:results} and~\ref{sec:generalization}. All configurations, greedy decoding included, are evaluated on all 15 test prompts.

\subsection{Inference}
\label{ssec:inference}
The unit of analysis for every confidence interval and hypothesis test is the \emph{prompt}. Two procedures were declared in advance:
\begin{itemize}[leftmargin=*]
  \item \textbf{Primary: exact sign-flip test.} For 15 prompts, all $2^{15}$ sign patterns are enumerated, giving an exact $p$-value for the paired difference in prompt-level RMST (\texttt{compute\_\allowbreak{}significance.py}, which reproduces every number in Table~\ref{tab:paired}).
  \item \textbf{Sensitivity: cluster bootstrap.} 3000 resamples of prompts give a 95\% CI for the same statistic.
\end{itemize}
At small $n$ the two need not agree; Section~\ref{ssec:disagreement} reports where they differ. The primary test governs.

\subsection{Metrics}
\textbf{RMST} has a horizon of 500 characters and uses Kaplan--Meier with right-censoring. Model-consistency NLL (mcNLL) scores every target position independently with its exact context window $\text{full}[\max(0,t{-}\text{block\_size}):t]$, batched with right-padding so that causal masking cannot leak filler into real positions (Appendix~\ref{app:defects}). 4-gram similarity is measured against a held-out reference corpus. For dual-calibrated configurations the exact $\KL(q\Vert p)$, achieved risk, and per-step feasibility diagnostics are reported.

%==============================================================================
\section{Method}
\label{sec:theory}
%==============================================================================

Everything in this section is local surrogate control in the sense of Section~\ref{sec:formal}: a one-step projection re-solved at every generated token from $\risk(v)$ alone. The theorem is exact, but it is a theorem about the one-step problem~(P), not about path-level control over $\tau$; Section~\ref{ssec:oracle} measures the distance between the two.

\subsection{The risk-bounded projection}
\begin{equation}
  \min_q \; \KL(q\Vert p) \quad \text{s.t.} \quad \E_q[\risk]\le\varepsilon,\quad q\in\Delta(\mathcal{V}). \tag{P}
\end{equation}
The Lagrangian gives, for fixed $\lambda\ge0$, $q_\lambda(v)\propto p(v)\exp(-\lambda\risk(v))$ (F). By weak duality this solves (P) \emph{only} for the $\lambda$ that satisfies complementary slackness ($\lambda{=}0$ if $\E_p[\risk]\le\varepsilon$, else the unique $\lambda$ with $\E_{q_\lambda}[\risk]=\varepsilon$); any other $\lambda$ solves the different, unconstrained problem $\min_q\KL(q\Vert p)+\lambda\E_q[\risk]$ (R). Three modes follow: \textsc{dual} solves (P) exactly, \textsc{fixed} solves (R), and \textsc{adaptive} is a heuristic with no guarantee.

\textbf{Edge cases.} If $\varepsilon\ge\E_p[\risk]$ the constraint is inactive and $\lambda^*{=}0$ exactly, so $q{=}p$ at zero distortion. If $\varepsilon<\min_v\risk(v)$, (P) is infeasible: even placing all mass on $\arg\min\risk$ gives $\E_q[\risk]=\min_v\risk(v)>\varepsilon$. At exactly $\varepsilon=\min_v\risk(v)$, (P) is feasible (the distribution concentrated on $\arg\min\risk$ attains equality) but no \emph{finite} $\lambda$ attains it through the exponential-tilting form (F): $g(\lambda):=\E_{q_\lambda}[\risk]$ approaches $\min_v\risk(v)$ only as $\lambda\to\infty$. In floating point this is usually invisible, since once non-minimal terms underflow to exactly $0.0$ the solver returns a correct feasible $q$. The comparison between a float32 \texttt{risk} tensor and a float64 \texttt{eps} can nevertheless round the wrong way at this point, so the solver uses an explicit tolerance and records a separate \texttt{near\_\allowbreak{}boundary} flag.

As $\lambda\to\infty$ for a single $n$-gram order, (F) converges to the same surviving set as hard no-repeat-$n$ blocking: both keep exactly the $\risk{=}0$ tokens, reweighted by $p$. The two differ in kind when no $\risk{=}0$ token exists. Hard blocking then bans the entire vocabulary, \texttt{torch.softmax} returns all NaN, and \texttt{torch.multinomial} raises a \texttt{RuntimeError}; (F) at a large finite $\lambda$ instead concentrates on the minimum-risk tokens. For the multi-order risk used here, $\lambda\to\infty$ concentrates on $\arg\min_v\risk(v)$, which matches the single-order limit only when $n_{\min}=n_{\max}$, so the method does not in general reduce to hard no-repeat blocking.

\subsection{Dual calibration}
\label{ssec:dual-fix}
The solver works entirely in log-space: $p$ is never materialized as a probability tensor; $q_\lambda$'s logits are $\log p-\lambda\risk$ computed from \texttt{F.log\_\allowbreak{}softmax} of the temperature-scaled logits, and $\KL(q\Vert p)$ and both entropies are computed from log-probabilities. This matters at low temperature: at $T{=}0.10$ many vocabulary entries have probabilities below the float32 range, and materializing $p$ would collapse distinct low-probability tokens to a single floor value.

Structural infeasibility is decided in closed form before any search: if $\min_v\risk(v)>\varepsilon$ (with tolerance), the step is reported as \texttt{structurally\_\allowbreak{}infeasible} and no search is attempted. Otherwise a feasible $\lambda$ exists, and the search bracket expands geometrically (doubling from $\lambda{=}1$, at most 60 doublings, so $\lambda\le2^{60}\approx1.15\times10^{18}$) until its upper end is feasible, and then bisects within it. This limit is practical and not mathematical; Appendix~\ref{app:lambda-estimate} shows that it is far beyond anything a risk-weighted logit shift requires away from the boundary case. Per-step diagnostics (feasibility status, $\lambda$, achieved risk, KL, and the boundary flag) are saved for every dual-mode configuration.

\subsection{Recurrence risk and prompt context}
$\risk(v)$ uses an incrementally maintained hash map, checked to be identical to brute-force history scanning. Whether the prompt's own tokens count as history was tested (\texttt{include\_\allowbreak{}prompt\_\allowbreak{}context}$\in\{$True,False$\}$): the RMST difference (492.9 vs.\ 499.9) is not significant under the primary exact test ($p{=}0.125$).

\subsection{The adaptive controller}
$\lambda_t=\lambda_{\text{base}}+\lambda_{\text{rep}}\max(0,\overline{\text{rep}}-\text{rep}_{\text{target}})-\lambda_{\text{ent}}\max(0,\overline{H}-H_{\text{target}})$, clipped to $[0,\lambda_{\max}]$. The lower clip is a correctness requirement, since an unclipped negative multiplier would reward risk.

\subsection{Baselines}
\textbf{FSD-style} (Section~\ref{sec:related}) samples from the probability-space contrast $p_\theta(v)-\alpha\,p_\omega(v)$ over the top-$k{=}6$ candidates, with $p_\omega$ the smoothed multi-order anti-LM built from the generated text. It uses a continuous, frequency-weighted contrast rather than a bounded projection with a distortion budget, and differs from repetition penalty in suppressing only candidates the anti-LM is confident about.

\textbf{LZ Penalty} \citep{ginart2025lzpenalty}. For a recent buffer $b$ (the last $B$ tokens) and an earlier lookback window $w$ (the $W$ tokens before it), each candidate $a$ receives a penalty from the longest match of the buffer extended by $a$ against the window: $\min(\log_2\delta_a,\log_2 V)-\log_2 V\in[-\log_2V,0]$, where $\delta_a$ is the distance to the closest window position achieving the longest match, and $0$ for a candidate with no match. This is added to the logits as $\alpha\cdot$penalty, so candidates completing a long, close repeat are suppressed most. Buffer and window sizes are scaled down from the paper's subword defaults (32/512) for the character-level vocabulary.

\textbf{SuffixMatch} finds the longest suffix of the generated history (last 400 tokens) that occurred earlier and subtracts $\alpha\cdot\text{match\_len}/20$ from the logits of the tokens that followed those earlier occurrences, with $\alpha{=}3$; it scans the history at every step, at cost $O(\text{history})$. All decoders share the same \texttt{step}/\texttt{reset}/\texttt{prime} interface and the same generation loop.


%==============================================================================
\section{Results}
\label{sec:results}
%==============================================================================

\subsection{Main comparison at $T{=}0.10$}
Table~\ref{tab:main} reports all configurations on the \textsc{cosine} checkpoint. Dual feasibility is exactly 100.0\% at both tested $\varepsilon$, with zero mean and zero maximum violation.

\begin{table}[h]
\centering
\caption{Primary (fiction-trained, \textsc{cosine}) checkpoint, $T{=}0.10$, all 15 test prompts, frozen event. \textsc{rr-dual} feasibility is 100.0\%.}
\label{tab:main}
\begin{tabular}{@{}lrrrrrr@{}}
\toprule
Method & $n$ & Loop \% & RMST & mcNLL & Sim & $\KL(q\Vert p)$ \\
\midrule
Greedy                      & 150 & 93.3\% & 130   & 0.820 & 0.089 & --- \\
Temperature 0.05             & 50  & 86.0\% & 170   & 0.831 & 0.116 & --- \\
Temperature 0.10             & 150 & 28.7\% & 404   & 0.886 & 0.155 & --- \\
Temperature 0.15             & 50  & 0.0\%  & 500   & 0.926 & 0.181 & --- \\
Repetition penalty 1.3        & 150 & 4.7\%  & 486   & 0.950 & 0.199 & --- \\
No-repeat 4-gram$^{*}$       & 150 & 0.0\%  & 500   & 1.722 & 0.249 & --- \\
SuffixMatch ($\alpha{=}3$)    & 150 & 0.0\%  & 500   & 1.017 & 0.267 & --- \\
FSD-style (reimpl.)          & 150 & 1.3\%  & 499   & 1.099 & 0.216 & --- \\
LZ Penalty (reimpl.)        & 150 & 0.0\%  & 500   & 0.924 & 0.180 & --- \\
\textsc{rr-fixed} ($\lambda{=}2$)   & 150 & 2.7\%  & 493   & 0.934 & 0.189 & 0.037 \\
\textsc{rr-adaptive}          & 150 & 0.0\%  & 500   & 0.960 & 0.212 & 0.096 \\
\textsc{rr-dual} ($\varepsilon{=}0.05$) & 150 & 0.0\% & 500 & 1.938 & 0.212 & 11.186 \\
\textsc{rr-dual} ($\varepsilon{=}0.01$) & 150 & 0.0\% & 500 & 2.054 & 0.193 & 12.670 \\
\bottomrule
\end{tabular}

\vspace{4pt}
{\footnotesize $^{*}$Hard constraint on the measured event; excluded from significance testing (Section~\ref{ssec:hard-repeat}).}
\end{table}

Distortion and residual risk are monotone across the three recurrence-risk modes: \textsc{fixed} (0.037 bits, 2.7\% loops), \textsc{adaptive} (0.096 bits, 0.0\%), \textsc{dual} (11.2--12.7 bits, 0.0\%). This describes the observed ordering of three points; no trend test was run. The dual mode's cost is what a properly feasible $\varepsilon$-bounded guarantee requires at this operating point, and the text it produces is visibly degraded (Section~\ref{ssec:qualitative}).

\subsection{No-repeat blocking is excluded from ranking}
\label{ssec:hard-repeat}
A global no-repeat-$n$-gram filter, checked against the entire history, structurally cannot allow a persistent cycle of period up to the generation length, since any such cycle reproduces a blocked $n$-gram at lag $P\le n$. Its 0/150 result holds by construction and is not an experimental finding, so it is excluded from every significance test and from the survival ranking.

\subsection{Paired comparisons}
\label{ssec:disagreement}
\begin{table}[h]
\centering
\caption{Paired, prompt-clustered comparisons of RMST (\texttt{compute\_\allowbreak{}significance.py}).}
\label{tab:paired}
\begin{tabular}{@{}lrrl@{}}
\toprule
Comparison ($A-B$) & $\Delta$RMST & $p$ (exact) & 95\% cluster CI \\
\midrule
\textsc{rr-fixed} $-$ Temperature 0.10   & $+88.4$  & $<0.001$ & [66.4, 113.9] \\
\textsc{rr-fixed} $-$ Rep.\ penalty 1.3  & $+6.5$   & $0.246$  & [$-3.7$, 15.9] \\
\textsc{rr-adaptive} $-$ Rep.\ penalty 1.3 & $+13.6$ & $0.016$ & [6.1, 21.9] \\
\textsc{rr-fixed} $-$ \{SuffixMatch, \textsc{dual}, \textsc{adaptive}, LZ Penalty\}$^{\dagger}$ & $-7.1$ & $0.125$ & [$-14.4$, $-1.1$] \\
\textsc{rr-fixed} $-$ FSD-style & $-5.9$ & $0.219$ & [$-13.9$, $0.9$] \\
\textsc{rr-adaptive} $-$ \textsc{rr-dual} & $0.0$ & $1.000$ & [0.0, 0.0] \\
FSD-style $-$ SuffixMatch & $-1.3$ & $0.500$ & [$-3.5$, $0.0$] \\
FSD-style $-$ LZ Penalty & $-1.3$ & $0.500$ & [$-3.5$, $0.0$] \\
LZ Penalty $-$ SuffixMatch$^{\dagger}$ & $0.0$ & $1.000$ & [0.0, 0.0] \\
\bottomrule
\end{tabular}

\vspace{4pt}
{\footnotesize $^{\dagger}$Identical because all listed configurations achieved 0/150; survival alone cannot distinguish them.}
\end{table}

The exact test and the bootstrap interval disagree on \textsc{rr-fixed} against the zero-loop group: the cluster interval excludes zero ($[-14.4,-1.1]$), but the pre-declared primary test gives $p{=}0.125$, not significant at 15 prompts. The primary test governs, and the difference is reported as not statistically significant. FSD-style's small survival deficit against the zero-loop methods (2/150 loops; $\Delta$RMST $-1.3$, $p{=}0.5$) is likewise not significant.

\subsection{LZ Penalty and FSD-style}
\label{ssec:results-fsd-lz}
LZ Penalty reaches 0/150 loops and joins SuffixMatch and \textsc{rr-adaptive} in the zero-loop group. FSD-style does not: it produces 2/150 loops (1.3\%) and is not a member of that group. Their spelling-error rates, computed from the generated text, are 2.59\% for FSD-style and 0.75\% for LZ Penalty, against 0.44\% for \textsc{rr-fixed} and 1.51\% for SuffixMatch. LZ Penalty's low error rate together with its low KL cost (Table~\ref{tab:commonkl}: 0.186 bits, against 2.184 bits for FSD-style) is consistent with its continuously varying, codelength-based penalty against FSD-style's coarser top-$k$ contrast, although no independent account of this difference is offered.

\subsection{Common-distortion comparison}
\label{ssec:common-kl}
Table~\ref{tab:main}'s KL column is populated only for the recurrence-risk modes, because only \texttt{RecurrenceRiskDecoder} computes a $q$ to compare against $p$; the other decoders sample from a well-defined $q$ at every step but carry no KL or risk figure. \texttt{common\_\allowbreak{}kl\_\allowbreak{}comparison.py} exposes every decoder's actual per-step sampling distribution (\texttt{last\_\allowbreak{}q}) and scores it against a shared reference $p$ (the temperature-scaled softmax of the raw model logits) and a shared standard recurrence-risk vector (computed by a passive shadow tracker that observes each decoder's generation history without influencing it).

\begin{table}[h]
\centering
\caption{Common-distortion comparison, \textsc{cosine} checkpoint, $T{=}0.10$, all 15 test prompts $\times$ 10 seeds ($n{=}150$ each). KL and risk against the SAME reference $p$ and SAME risk definition for every row -- not each decoder's own bookkeeping. Sorted by KL.}
\label{tab:commonkl}
\begin{tabular}{@{}lrrrr@{}}
\toprule
Method & $\KL(q\Vert p)$ (bits) & $\E_q[\risk]$ & Latency (ms/tok) & Loop \% \\
\midrule
Temperature 0.10 (reference)   & 0.000 & 0.686 & 1.57 & 28.7\% \\
\textsc{rr-fixed}               & 0.037 & 0.550 & 3.05 & 2.7\% \\
\textsc{rr-adaptive}            & 0.096 & 0.466 & 3.18 & 0.0\% \\
LZ Penalty                      & 0.186 & 0.624 & 7.91 & 0.0\% \\
Repetition penalty 1.3          & 0.388 & 0.533 & 7.88 & 4.7\% \\
SuffixMatch ($\alpha{=}3$)      & 0.436 & 0.285 & 1.63 & 0.0\% \\
FSD-style                       & 2.184 & 0.405 & 12.92 & 1.3\% \\
No-repeat 4-gram$^{*}$         & 8.796 & 0.066 & 1.64 & 0.0\% \\
\textsc{rr-dual} ($\varepsilon{=}0.05$) & 11.186 & 0.017 & 6.49 & 0.0\% \\
\textsc{rr-dual} ($\varepsilon{=}0.01$) & 12.670 & 0.003 & 6.54 & 0.0\% \\
\bottomrule
\end{tabular}

\vspace{4pt}
{\footnotesize $^{*}$Hard constraint on the measured event, as elsewhere. 95\% cluster CIs for KL and risk are within $\pm$6\% of the mean shown; latency CIs are wider for FSD-style ([8.5, 16.9] ms/token) and LZ Penalty ([5.9, 10.2]). See \texttt{runs/common\_\allowbreak{}kl\_\allowbreak{}comparison\_\allowbreak{}v2\_\allowbreak{}merged/common\_\allowbreak{}kl\_\allowbreak{}comparison.csv}.}
\end{table}

The four recurrence-risk rows reproduce Table~\ref{tab:main} to three decimals (0.037, 0.096, 11.186, 12.670 bits) through an independent code path (\texttt{last\_\allowbreak{}q} and \texttt{safe\_\allowbreak{}kl\_\allowbreak{}bits} rather than \texttt{RecurrenceRiskDecoder.diagnostics()}), a consistency check that found no discrepancy.

\textbf{Cost of zero loops.} \textsc{rr-adaptive} is the cheapest route to zero measured loops overall (0.096 bits), but carries no formal guarantee. Among the external baselines that reach exactly zero loops (LZ Penalty, SuffixMatch, no-repeat 4-gram), LZ Penalty is cheapest (0.186 bits). FSD-style reaches 1.3\% loops at 2.184 bits, an order of magnitude costlier than LZ Penalty and not zero-loop. Repetition penalty is comparably cheap in KL (0.388 bits) but only partially suppresses the event (4.7\% loops). Together with FSD-style it is also among the slowest methods tested (7.88 and 12.92 ms/token, roughly $4$--$8\times$ the exact-match baselines), since both rescan or re-score the recent history at every step. No-repeat 4-gram, often treated as free because it spends no explicit budget, has the second-highest KL measured (8.796 bits): a hard mask forces exactly the redistribution that the dual solver's boundary case describes, without reporting its cost.

\textbf{Risk and KL do not lie on a single frontier across families.} LZ Penalty reaches lower KL than SuffixMatch (0.186 vs.\ 0.436 bits) but higher achieved risk (0.624 vs.\ 0.285): neither dominates the other on this pair of axes, though both reach 0/150 loops. Within a single family the expected monotone relationship holds (\textsc{rr-fixed}$\to$\textsc{rr-adaptive}$\to$\textsc{rr-dual} $\varepsilon{=}0.05\to\varepsilon{=}0.01$ strictly lowers achieved risk as KL rises), but it does not extend across families, which is the reason for measuring all of them against one reference.

\subsection{Qualitative example}
\label{ssec:qualitative}
Same prompt, same checkpoint, $T{=}0.10$:
\begin{quote}\small
\textbf{Greedy} (loops): \emph{``they were all the subject of the consideration of the subject of the subject of the subject of the subject of the subject of the subject of the subject of the subject of the consideration of the subject of the subject of''}

\textbf{\textsc{rr-dual} $\varepsilon{=}0.05$} (no cycle, visibly degraded): \emph{``they had been a sensible of tranquillity and misery to take at time in terms on their sistence was always assured, after two suppose to me, with a considerably, tlength t\^ete hear or tyle so are,--''}

\textbf{LZ Penalty} (no cycle, more coherent than dual): \emph{``they were all the conceiving which I had been so much as the stranger, and the proposal of the contrary of the present which I was always to be done the most and suffering the strong of the conversation which I had been property to the most present and all the''}
\end{quote}
The dual example is the most degraded of the three, consistent with its KL cost being the highest (11.2--12.7 bits). The blinded human evaluation (Section~\ref{sec:human-eval}) tests whether this impression holds across readers.

%==============================================================================
\section{Human Evaluation}
\label{sec:human-eval}
%==============================================================================

Task 1 reports participant agreement with the frozen event's classification, not with ground truth. The event (Section~\ref{sec:event}) is a calibrated operational detector with a known, nonzero false-positive rate concentrated on specific corpora (Section~\ref{ssec:discussion}); Task 1 measures whether human judgment tracks this specific syntactic definition, which already disagrees with a semantic reading on borderline cases.

A blinded, self-contained offline page (\texttt{human\_\allowbreak{}eval\_\allowbreak{}tool.html}, with an English or Russian interface) has two tasks: (1) loop-detection agreement on 18 passages (clearly looping, clearly non-looping, and borderline), and (2) blind pairwise preference on 12 pairs covering all six pairings among the four methods that reach 0/150 on the primary model (\textsc{dual}, SuffixMatch, FSD-style, LZ Penalty). Automatic metrics alone cannot rank this group: SuffixMatch's lower mcNLL is not a validated quality signal, and FSD-style and LZ Penalty's spelling-error rates point in different directions from their mcNLL. The tool randomizes presentation order and left/right side independently per participant and ends with a copy-to-clipboard export of the raw responses (\texttt{aggregate\_\allowbreak{}human\_\allowbreak{}eval.py} combines the exports).

Nine participants were recruited informally, not through a formal panel, and sessions took a median of about 9.6 minutes (range 2--11 minutes). Each raw export was checked for internal consistency and independence before inclusion, and only sessions that passed are reported.

\begin{table}[h]
\centering
\caption{Human evaluation, 9 participants. Task 1: loop-detection agreement against the frozen event's own classification (an operational instrument, not ground truth -- see note above). Task 2: blind pairwise preference, all six pairings among the zero-loop group; ``first'' is the method named first in each row.}
\label{tab:humaneval}
\begin{tabular}{@{}lrl@{}}
\toprule
Task 1 (loop detection) & & \\
\midrule
Pooled accuracy & 140/162 & 86.4\%, 95\% CI [78.4\%, 93.2\%] \\
\quad Looping items & 54/54 & 100.0\% \\
\quad Non-looping items & 54/54 & 100.0\% \\
\quad Borderline items & 32/54 & 59.3\% \\
Inter-participant agreement & 540/648 & 83.3\% \\
\midrule
Task 2 (pairwise preference, by participant) & & \\
\midrule
\textsc{dual} vs.\ FSD-style & 2/7 & \textsc{dual} 29\%, 95\% CI [4\%, 71\%] \\
\textsc{dual} vs.\ LZ Penalty & 3/5 & \textsc{dual} 60\%, 95\% CI [15\%, 95\%] \\
\textsc{dual} vs.\ SuffixMatch & 3/7 & \textsc{dual} 43\%, 95\% CI [10\%, 82\%] \\
FSD-style vs.\ LZ Penalty & 2/3 & FSD-style 67\%, 95\% CI [9\%, 99\%] \\
FSD-style vs.\ SuffixMatch & 2/4 & FSD-style 50\%, 95\% CI [7\%, 93\%] \\
LZ Penalty vs.\ SuffixMatch & 1/3 & LZ Penalty 33\%, 95\% CI [1\%, 91\%] \\
\bottomrule
\end{tabular}

\vspace{4pt}
{\footnotesize Task 1's CI is a participant-clustered bootstrap (10,000 resamples of the 9 participants, not the 162 individual judgments): pooling repeated judgments from the same rater as if they were independent trials (plain Clopper-Pearson) understates uncertainty, since a systematically stricter or looser rater pulls all of their own judgments the same direction. Task 2 is reported per participant (each participant's majority preference across their 2 votes per pairing collapsed to one data point, participants who split their votes evenly or tied on both excluded, hence $n\le9$ per row) with an exact (Clopper-Pearson) binomial CI, which remains valid at small $n$ and when all participants agree. See \texttt{aggregate\_\allowbreak{}human\_\allowbreak{}eval.py}'s \texttt{cluster\_\allowbreak{}level\_\allowbreak{}clopper\_\allowbreak{}pearson} (Task 2) and \texttt{participant\_\allowbreak{}cluster\_\allowbreak{}bootstrap\_\allowbreak{}ci} (Task 1).}
\end{table}

\textbf{Loop detection is reliable on clear cases and imperfect on the syntactic event's blind spot.} Agreement with the frozen event's classification is 100.0\% on both unambiguously looping and unambiguously non-looping items and 59.3\% on the borderline items, which are included deliberately because they probe the syntactic-not-semantic nature of the event (Section~\ref{ssec:discussion}). All 22 pooled disagreements are false positives (participants calling borderline non-loops ``looping'') and none are false negatives, consistent with readers being more willing than the syntactic event to call something a loop.

\textbf{No zero-loop decoder is clearly preferred.} Across the 108 pairwise judgments (27 ties), pooled win rates are 43\% for FSD-style and SuffixMatch, 35\% for \textsc{dual}, and 30\% for LZ Penalty, and every participant-level 95\% interval in Table~\ref{tab:humaneval} includes 50\%. Individual participants differ strongly: one preferred \textsc{dual} in 6 of 12 comparisons, another called 10 of 12 pairs a tie. The data neither support nor contradict a readability cost for \textsc{dual}. Its KL cost (11.2--12.7 bits, Section~\ref{ssec:common-kl}) is far larger than that of the heuristic decoders that reach the same observed 0/150 outcome, but at this sample size that cost is not matched by a detectable difference in human-judged readability, and the automatic proxies (mcNLL, spelling-error rate) should not be read as predicting the human judgments collected here.

This is a pilot at a single temperature ($T{=}0.10$) on the primary model with nine participants; a properly powered study with a larger, more systematically recruited pool is a natural next step.

%==============================================================================
\section{Toward Finite-Horizon Hazard: Synthetic Validation}
\label{sec:synthetic}
%==============================================================================

\subsection{Construction}
A 20-state Markov chain has a diffuse ``safe'' region from which a small probability (0.05) leads to a \emph{gateway} state, which is not itself an immediate repeat but transitions with high probability (0.75) into a tight 2-cycle. The true hazard is computed by exact backward dynamic programming over (state, bounded 2-state trailing window, steps remaining): $f[r][(s,\text{win})]=\sum_{s'}P[s,s']\big(\mathbb{1}[s'\in\text{win}]+\mathbb{1}[s'\notin\text{win}]\,f[r{-}1][(s',\text{win}')]\big)$, with $f[0][\cdot]{=}0$. The per-candidate hazard $\haz(s')$ is the probability of looping within the remaining horizon given that $s'$ is chosen now.

\subsection{Matched-distortion comparison}
Hazard-aware control's one-step budget is bisected exactly, without Monte Carlo, until its own sequence-level loop probability equals $\varepsilon$, so ``Hazard $\varepsilon$'' and ``Hazard loop\%'' in Table~\ref{tab:synthetic} agree by construction. The matched-distortion comparison is computed and saved for every point of the $\varepsilon$ grid at once, so no operating point is hand-picked.

\begin{table}[h]
\centering
\caption{Matched-distortion comparison, computed for every point of the $\varepsilon$ grid.}
\label{tab:synthetic}
\begin{tabular}{@{}rrrrrll@{}}
\toprule
Hazard $\varepsilon$ & Hazard KL & Hazard loop\% & Matched local $\varepsilon$ & Local KL & Local loop\% & Comparable? \\
\midrule
0.30 & 0.255 & 30.0\% & 0.263 & 0.255 & 43.9\% & yes \\
0.20 & 0.245 & 20.0\% & 0.280 & 0.248 & 44.2\% & yes \\
0.10 & 0.293 & 10.0\% & 0.196 & 0.291 & 41.1\% & yes \\
0.05 & 0.412 & 5.0\%  & 0.052 & 0.415 & 24.9\% & yes \\
0.02 & 0.533 & 2.0\%  & 0.018 & 0.511 & 9.6\%  & yes \\
0.01 & 0.632 & 1.0\%  & 0.001 & 0.602 & 0.6\%  & \textbf{no} \\
\bottomrule
\end{tabular}
\end{table}

\textbf{KL convention.} Every KL figure in this section is \emph{absorbing} at the loop event: a trajectory that loops at step $t$ pays its accumulated one-step distortion up to and including the decision that caused the loop, and nothing afterward. This matches the oracle's exact derivation below, in which the backward recursion's ``loop happens now'' branch terminates with a fixed weight.

Local recurrence-risk control has an intrinsic per-step KL ceiling that hazard-aware control does not share. Sweeping $\varepsilon\to0$ directly, local's achievable KL saturates at $\approx0.613$ bits/step and does not increase below $\varepsilon{=}0.0001$, a structural limit and not a scan-resolution artifact. It is the same phenomenon as the boundary $\varepsilon=\min_v\risk(v)$ in the real dual solver (Section~\ref{sec:theory}): once $q$ is fully concentrated on the one-step risk-minimizing actions, tightening $\varepsilon$ cannot change $q$. Hazard's KL at $\varepsilon{=}0.01$ (0.632) exceeds this ceiling, so no local configuration can match its distortion there; that point is marked not comparable and excluded from the dominance count. Among the five points where a matched-KL comparison is possible ($\varepsilon\ge0.02$), hazard-aware control has the lower loop rate at all five, by margins of 7.6--31.1 percentage points. At $\varepsilon{=}0.20$ the exact values are 20.0\% (hazard) against 44.2\% (local).

\subsection{The path-level event-control oracle: how inefficient is local surrogate control?}
\label{ssec:oracle}
Local and hazard-aware control are both local surrogate control: one-step projections that re-solve independently at every step from whatever risk signal they use (an indicator for local, the exact multi-step hazard for hazard-aware), and neither accounts for the KL cost of the choices it will face later. The true path-level optimum is the exact path measure $Q^*$ minimizing total $\KL(Q\Vert P)$ subject to $Q(\text{loop within }H)\le\varepsilon$, that is, $P_Q(\tau\le H)\le\varepsilon$ in the notation of Section~\ref{sec:formal}. It is given by the standard exponential-tilting (Doob-transform) solution of a KL-constrained hitting-time problem; the construction is not novel, and its purpose is to measure what local surrogate control gives up relative to it. Let $g[r][(s,\text{win});\lambda]:=\E_P[2^{-\lambda\cdot\mathbb{1}(\text{loop within }r)}\mid s,\text{win}]$. It satisfies the same backward recursion as $f$, with the ``loop happens now'' branch weighted by $2^{-\lambda}$ instead of probability mass $1$. The optimal tilted kernel is $Q^*(s'\mid s,\text{win},r)\propto P(s,s')\cdot m(s',\text{win}',r{-}1)$ (the per-branch term defining $g$), and $\lambda$ is bisected, as in the per-step dual solver but over the whole horizon, until $Q^*$'s exactly forward-computed loop probability matches each target $\varepsilon$.

The oracle's $\varepsilon$ bounds the \emph{sequence-level} probability $Q(\text{loop within }H)$, whereas the one-step solvers' own $\varepsilon$ bounds the one-step expected risk $\E_q[\risk]$. To compare all three mechanisms at the same operating point, local and hazard-aware control's one-step budgets are bisected (exactly, by the same finite-state forward propagation, without Monte Carlo) until their own sequence-level loop probability also matches the target. The one-step solver's feasibility flag is tracked and reported (Infeasible column): hazard's one-step bound is infeasible at 33--50\% of (probability-mass-weighted) steps across the grid, whereas local is never infeasible, so hazard's sequence-level target is met only because unconstrained steps compensate for steps pinned at their achievable floor.

\begin{table}[h]
\centering
\caption{Sequence-level oracle vs.\ local and hazard-aware one-step control, all matched to the same sequence-level $\varepsilon$. The oracle's achieved $\varepsilon$ matches the target to numerical precision (residual $<10^{-10}$), and local and hazard-aware achieved $\varepsilon$ match it to within bisection tolerance ($<10^{-9}$); all are exact, not simulated. \textbf{Infeasible}: probability-mass-weighted fraction of steps at which the one-step budget was structurally below hazard's achievable floor.}
\label{tab:oracle}
\begin{tabular}{@{}rrrrrrrr@{}}
\toprule
$\varepsilon$ & Local KL/step & Hazard KL/step & \textbf{Oracle KL/step} & Local/Oracle & Hazard/Oracle & Hazard Infeasible \\
\midrule
0.30 & 0.390 & 0.255 & \textbf{0.077} & 5.1$\times$ & 3.3$\times$ & 33\% \\
0.20 & 0.443 & 0.245 & \textbf{0.123} & 3.6$\times$ & 2.0$\times$ & 50\% \\
0.10 & 0.511 & 0.293 & \textbf{0.184} & 2.8$\times$ & 1.6$\times$ & 50\% \\
0.05 & 0.555 & 0.412 & \textbf{0.224} & 2.5$\times$ & 1.8$\times$ & 50\% \\
0.02 & 0.587 & 0.533 & \textbf{0.254} & 2.3$\times$ & 2.1$\times$ & 50\% \\
0.01 & 0.599 & 0.632 & \textbf{0.266} & 2.3$\times$ & 2.4$\times$ & 50\% \\
\bottomrule
\end{tabular}
\end{table}

\textbf{Even hazard-aware control, using the true multi-step hazard as its risk signal, spends 1.6--3.3$\times$ the oracle's KL (mean 2.2$\times$) for the same loop-suppression level; local spends 2.3--5.1$\times$ (mean 3.1$\times$).} Perfect knowledge of future risk does not close the gap to the sequence-level optimum, because a one-step projection cannot account for the distortion cost of the choices it will be forced into later. Hazard is cheaper than local at 5 of 6 grid points, by a margin that widens as $\varepsilon$ loosens; at the tightest budget ($\varepsilon{=}0.01$) local is slightly cheaper (2.3$\times$ against 2.4$\times$), where hazard's structural infeasibility binds most.

\textbf{Why an exact hazard signal is not enough.} The exact one-step marginal of $Q^*$ has a closed form that differs from the exponential-tilting form hazard-aware control uses. Write $H(h,v)$ for the exact hazard (the probability of looping within the remaining horizon given that candidate $v$ is chosen from history $h$, with $H(h,v){:=}1$ if $v$ completes an immediate loop). Because ``loop within $r$ steps'' is a single Bernoulli event, $g[r]=\E_P[2^{-\lambda\cdot\mathbb{1}(L)}]=(1-f[r])+f[r]\,2^{-\lambda}$, that is,
\[ g[r][(s,\text{win})]=1-(1-2^{-\lambda})\,f[r][(s,\text{win})], \]
for every $(r,s,\text{win})$ and a single $\lambda$ shared across the horizon; this was verified numerically to $<10^{-14}$ across all 58{,}940 nodes of the problem's $(r,s,\text{win})$ table. Substituting into $Q^*$'s kernel gives
\[ Q^*(v\mid h)\;\propto\;P(v\mid h)\cdot\bigl[\,1-(1-2^{-\lambda})\,H(h,v)\,\bigr], \]
with one globally calibrated $\lambda$ (a single bisection on the sequence-level target, not re-solved per step). Hazard-aware control instead uses $q(v\mid h)\propto P(v\mid h)\cdot2^{-\lambda\cdot H(h,v)}$, the same exact signal $H$ reweighted by the wrong functional form: exponential in hazard rather than the optimal affine form. The shortfall of hazard-aware control is therefore not only that it looks too little ahead (it uses exactly the information $H$ gives it) but that it converts exact knowledge of future risk into a token distribution in a suboptimal way. Whether the closed form suggests a tractable mechanism for a real, exponentially large state space is left open: $H$ needs an estimator there, and Section~\ref{sec:hazard-transfer} reports that the rollout estimator tried does not yield a usable one.

%==============================================================================
\section{Attempted Transfer to the Real Model}
\label{sec:hazard-transfer}
%==============================================================================

\subsection{Design}
Per-candidate hazard labels were collected by forcing each of several candidate next tokens (the model's top 8 by probability plus 4 lower-probability tokens) as the immediate next token, then continuing $K$ independent stochastic rollouts at $T{=}0.10$ for a fixed horizon and checking the frozen loop event. The horizon and context selection were designed so that the event can be observed: contexts are sampled near actual loops rather than uniformly at random.

\subsection{No detectable signal under this estimator}
Per-candidate hazard for ``risky'' ($\risk(v)>0$) and ``novel'' ($\risk(v)=0$) candidates shows no separation (0.027 vs.\ 0.028 mean hazard, $n{=}156$ and $1644$). A multi-step-commitment variant shows no significant separation at any tested commitment length $L\in\{1,3,5,8,15\}$ (largest raw gap $p{=}0.137$ at $L{=}8$; Bonferroni-corrected $p{=}0.685$).

The conclusion is deliberately narrow: under this rollout-based estimator, candidate-selection scheme, and per-cell sample size (5 rollouts per candidate), no hazard signal distinguishing risky from novel candidates was detected at single-token or multi-step-commitment granularity. This does not show that finite-horizon hazard control fails on real language models, that the mechanism of Section~\ref{sec:synthetic} is invalid, or that no estimator could find a signal. A properly powered, context-clustered re-estimation with many more rollouts per candidate is a natural next step.

%==============================================================================
\section{Generalization}
\label{sec:generalization}
%==============================================================================

\subsection{Pretrained subword pilot}
Before a full comparison on a pretrained subword model, a cheap pilot checks that a measurable low-temperature repetition regime exists, without assuming that the character model's temperature range or the pathology itself transfers. GPT-2 (124M) with an uncalibrated placeholder token-level event ($R{=}5$, $P_{\max}{=}20$ tokens) gives Table~\ref{tab:gpt2pilot}.

\begin{table}[h]
\centering
\caption{GPT-2 pilot, 5 prompts $\times$ 5 seeds, 150 new tokens, uncalibrated placeholder event.}
\label{tab:gpt2pilot}
\begin{tabular}{@{}rrl@{}}
\toprule
$T$ & Loop rate & 95\% CI \\
\midrule
0.05 & 32.0\% & [14.9\%, 53.5\%] \\
0.10 & 32.0\% & [14.9\%, 53.5\%] \\
0.15 & 40.0\% & [21.1\%, 61.3\%] \\
0.20 & 28.0\% & [12.1\%, 49.4\%] \\
0.30 & 28.0\% & [12.1\%, 49.4\%] \\
0.50 & 4.0\%  & [0.1\%, 20.4\%] \\
0.70 & 0.0\%  & [0.0\%, 13.7\%] \\
1.00 & 0.0\%  & [0.0\%, 13.7\%] \\
\bottomrule
\end{tabular}
\end{table}

A measurable regime exists (peak 40.0\% at $T{=}0.15$, above the pilot's 10\% threshold), so the full comparison was run (Section~\ref{sec:gpt2-full}; see Section~\ref{ssec:inference} for the inference protocol) after replacing the placeholder with a calibrated event, following the calibrate/freeze/validate protocol of Section~\ref{sec:event} at the token level on this model's own tokenizer. The frozen token-level event ($R{=}5$, $P_{\max}{=}20$ tokens; \texttt{validate\_\allowbreak{}loop\_\allowbreak{}event\_\allowbreak{}subword.py}) has a pooled false-positive rate of 1.1\% (95\% CI [0.7\%, 1.9\%]) on a disjoint validation battery. It coincides numerically with the pilot's placeholder, but was calibrated and validated before use.

\subsection{GPT-2 full comparison}
\label{sec:gpt2-full}
Table~\ref{tab:gpt2full} reports the same decoder classes as Table~\ref{tab:main}, run unmodified against GPT-2 through a thin HuggingFace generation loop (\texttt{gpt2\_\allowbreak{}full\_\allowbreak{}comparison.py}) that delegates every step to the same \texttt{.step()} interface. $T{=}0.10$; 10 prompts (written fresh and not inspected while choosing $T$, $\alpha$, or $\varepsilon$) $\times$ 10 seeds $\times$ 200 new tokens; risk order $n_{\min}{=}2,n_{\max}{=}4$ tokens (a phrase-length band appropriate to BPE tokens); LZ Penalty at the paper's subword defaults (buffer 32, window 512).

\begin{table}[h]
\centering
\caption{GPT-2 (124M), $T{=}0.10$, 10 prompts $\times$ 10 seeds, 200 new tokens, frozen validated event ($R{=}5$, $P_{\max}{=}20$ tokens).}
\label{tab:gpt2full}
\begin{tabular}{@{}lrrrr@{}}
\toprule
Method & $n$ & Loop \% & RMST (tok) & mcNLL \\
\midrule
Greedy                      & 100 & 30.0\% & 158.5 & 0.626 \\
Temperature 0.05             & 100 & 30.0\% & 167.2 & 0.678 \\
Temperature 0.10             & 100 & 34.0\% & 163.7 & 0.728 \\
Temperature 0.15             & 100 & 34.0\% & 164.1 & 0.766 \\
Repetition penalty 1.3        & 100 & 0.0\%  & 200.0 & 4.282 \\
No-repeat 3-gram$^{*}$       & 100 & 0.0\%  & 200.0 & 2.347 \\
SuffixMatch ($\alpha{=}3$)    & 100 & 4.0\%  & 198.3 & 1.203 \\
FSD-style (reimpl.)          & 100 & 0.0\%  & 200.0 & 2.122 \\
LZ Penalty (reimpl.)        & 100 & 8.0\%  & 193.6 & 1.532 \\
\textsc{rr-fixed} ($\lambda{=}3$)   & 100 & 22.0\% & 180.7 & 0.847 \\
\textsc{rr-adaptive}          & 100 & 19.0\% & 187.1 & 0.884 \\
\textsc{rr-dual} ($\varepsilon{=}0.05$) & 100 & 0.0\% & 200.0 & 2.525 \\
\textsc{rr-dual} ($\varepsilon{=}0.01$) & 100 & 0.0\% & 200.0 & 2.553 \\
\bottomrule
\end{tabular}

\vspace{4pt}
{\footnotesize $^{*}$Hard constraint on the measured event; excluded from the pattern discussed below on the same grounds as Table~\ref{tab:main}.}
\end{table}

Both dual configurations reach exactly 100.0\% per-step feasibility and 0/100 loops at 6.349--6.721 bits of KL: the same outcome as the primary model at roughly half its KL cost (11.2--12.7 bits), a difference that may reflect what concentrating away from a handful of risky tokens costs in a much larger subword vocabulary. No common numeric constant is claimed; only the dual mechanism's per-step feasibility, which does not depend on a corpus- or model-tuned constant, again reaches its target exactly.

Configurations without the dual guarantee do not carry their primary-model effectiveness over. \textsc{rr-fixed} and \textsc{rr-adaptive}, essentially loop-free on the primary model (2.7\%, 0.0\%), reach only 22.0\% and 19.0\% loops on GPT-2 at the same hyperparameters. LZ Penalty rises from 0.0\% on the primary model to 8.0\% here, and SuffixMatch is close to zero but not exactly zero (4.0\%). FSD-style moves in the opposite direction from its primary-model behavior (1.3\% there, 0.0\% here), so being close to zero on one model does not predict which way a heuristic's residual rate moves on another. Repetition penalty and hard no-repeat fully suppress the event here (0.0\% each), more effectively than on the primary model (4.7\% for repetition penalty), at a steep mcNLL cost (4.282 and 2.347).

\subsection{A second, independent character-model setting}
To test whether the primary result is specific to one 848K-parameter model on one corpus, a second model of identical architecture was trained on a corpus disjoint in genre, register, and author from every corpus used elsewhere: Darwin's \emph{On the Origin of Species} and Adam Smith's \emph{The Wealth of Nations} for training, Thoreau's \emph{Walden} for validation, and Machiavelli's \emph{The Prince} for test (non-fiction, argumentative prose, against the primary model's dialogue- and narrative-heavy fiction). The loop event was re-confirmed and not re-tuned on this corpus's held-out text (a \texttt{--confirm\_\allowbreak{}only} mode reusing the frozen $R{=}5$, $P_{\max}{=}60$ setting, which avoids calibrating and validating on the same corpus): pooled false-positive rate 5/600 (0.8\%, 95\% CI [0.3\%, 1.9\%]), comparable to the primary battery's $\sim$1.3\%, and the synthetic sensitivity controls again behaved correctly on 98--100\% of constructed cases. Table~\ref{tab:generalization} gives the comparison.

\begin{table}[h]
\centering
\caption{Second, non-fiction-trained model, $T{=}0.10$, all 15 test prompts, re-confirmed event. Hyperparameters unchanged from Table~\ref{tab:main}.}
\label{tab:generalization}
\begin{tabular}{@{}lrrrrr@{}}
\toprule
Method & $n$ & Loop \% & RMST & mcNLL & $\KL(q\Vert p)$ \\
\midrule
Greedy                      & 150 & 100.0\% & 269   & 0.593 & --- \\
Temperature 0.05             & 50  & 100.0\% & 325   & 0.603 & --- \\
Temperature 0.10             & 150 & 64.7\%  & 396   & 0.617 & --- \\
Temperature 0.15             & 50  & 26.0\%  & 458   & 0.641 & --- \\
Repetition penalty 1.3        & 150 & 36.0\%  & 458   & 0.648 & --- \\
No-repeat 4-gram$^{*}$       & 150 & 0.0\%   & 500   & 1.605 & --- \\
SuffixMatch ($\alpha{=}3$)    & 150 & 0.0\%   & 500   & 0.706 & --- \\
FSD-style (reimpl.)          & 150 & 16.7\%  & 486   & 0.852 & --- \\
LZ Penalty (reimpl.)        & 150 & 2.0\%   & 498   & 0.649 & --- \\
\textsc{rr-fixed} ($\lambda{=}2$)   & 150 & 51.3\%  & 439   & 0.615 & 0.018 \\
\textsc{rr-adaptive}          & 150 & 21.3\%  & 479   & 0.633 & 0.077 \\
\textsc{rr-dual} ($\varepsilon{=}0.05$) & 150 & 0.0\% & 500 & 1.994 & 12.891 \\
\textsc{rr-dual} ($\varepsilon{=}0.01$) & 150 & 0.0\% & 500 & 2.146 & 14.727 \\
\bottomrule
\end{tabular}
\end{table}
{\footnotesize $^{*}$Hard constraint on the measured event; excluded from significance testing, as in Table~\ref{tab:main}.}

\subsection{What transfers and what does not}
\textbf{The per-step feasibility rate is consistent across models; zero persistent loops is an empirical observation, not a proven consequence.} \textsc{rr-dual} reaches exactly 100.0\% per-step feasibility on this corpus at both $\varepsilon$, identical to the primary model. Feasibility is the quantity the theorem governs: it bounds the per-step \emph{expected} recurrence risk when the numerical projection succeeds, and its success rate stays consistent across corpora because $\lambda$ is re-solved from $p$ and $\risk$ at every step, with no corpus-tuned constant. The theorem says nothing about the sequence-level count of persistent loops, so 0/150 on both models is reported as a repeated empirical observation.

\textbf{The fixed-hyperparameter modes do not transfer.} \textsc{rr-fixed} ($\lambda{=}2$, unchanged) rises from 2.7\% loops on the primary model to 51.3\% here, and \textsc{rr-adaptive} from 0.0\% to 21.3\%. The gap between \textsc{rr-dual} and \textsc{rr-fixed} on this corpus is large and highly significant (paired exact test, $\Delta$RMST $=+60.8$, $p{=}0.0001$). A fixed penalty strength suited to one corpus's risk and entropy statistics does not carry its effectiveness to a corpus with different repetition statistics, although the code and hyperparameters are identical. \textsc{rr-fixed}'s output here remains visibly cyclical (\emph{``the same manner as the produce of the country, the country which they have been sufficient to the produce of the same manner, and the profits of the same manner as the produce of the produce of the land and labour of the produce of the produce of''}).

\textbf{The primary model's temperature-safe regime does not transfer.} $T{=}0.15$, which reached 0.0\% loops on the primary model, still shows 26.0\% here. The temperature sweep was not repeated for this model, so its safe threshold is unknown.

Repetition penalty and \textsc{rr-adaptive} weaken substantially on this corpus (36.0\% and 21.3\% loops, from 4.7\% and 0.0\%). Among the history-aware and hard-constraint mechanisms, effectiveness transfers unevenly: SuffixMatch and hard no-repeat still reach 0/150, whereas FSD-style rises to 16.7\% (from 1.3\%) and LZ Penalty to 2.0\% (from 0.0\%). LZ Penalty's effectiveness also does not carry over to GPT-2 (0.0\% to 8.0\%). For FSD-style and LZ Penalty, unlike SuffixMatch and hard no-repeat, effectiveness is therefore a property of the model and corpus and not of the mechanism's design.

\subsection{A compact validation across modern pretrained model families}
\label{ssec:modern-llm}
To test whether any of this transfers to modern pretrained models beyond GPT-2, three checkpoints spanning roughly an order of magnitude in size and three pretraining organizations were selected: OLMo-2-0425-1B (Allen AI, 1B), Qwen2.5-3B (Alibaba, 3B), and Mistral-7B-v0.3 (Mistral AI, 7B). Each follows the pilot, calibrate, compare protocol used for GPT-2, at the same compact scale across all three: 5 seeds $\times$ 10 fresh prompts $\times$ 150 new tokens (50 samples per configuration).

All three models run in float32, matching the primary and GPT-2 experiments and not the faster bfloat16 default of this project's CUDA code, so that numerical precision does not vary across the models being compared. Mistral-7B-v0.3's weights do not fit in the 36GB unified-memory hardware under the default HuggingFace loading path, so they were loaded shard by shard directly onto the target device (\texttt{device\_\allowbreak{}map} with \texttt{low\_\allowbreak{}cpu\_\allowbreak{}mem\_\allowbreak{}usage=True}), verified bit-identical to the default path on a smaller model.

\begin{table}[h]
\centering
\caption{Pilot: loop rate (\%) by temperature, 5 prompts $\times$ 5 seeds, 150 new tokens, uncalibrated placeholder event ($R{=}5$, $P_{\max}{=}20$ tokens) per model's own tokenizer.}
\label{tab:modernpilot}
\begin{tabular}{@{}rrrr@{}}
\toprule
$T$ & OLMo-2-0425-1B & Mistral-7B-v0.3 & Qwen2.5-3B \\
\midrule
0.05 & 32.0\% & 12.0\% & 0.0\% \\
0.10 & \textbf{44.0\%} & \textbf{20.0\%} & 0.0\% \\
0.15 & 16.0\% & 8.0\%  & 0.0\% \\
0.20 & 16.0\% & 12.0\% & 0.0\% \\
0.30 & 12.0\% & 8.0\%  & 0.0\% \\
0.50 & 8.0\%  & 0.0\%  & 0.0\% \\
0.70 & 0.0\%  & 0.0\%  & 0.0\% \\
1.00 & 0.0\%  & 0.0\%  & 0.0\% \\
\bottomrule
\end{tabular}
\end{table}

\textbf{Qwen2.5-3B shows no measurable regime at any tested temperature} (Table~\ref{tab:modernpilot}), 0/25 at each of eight temperatures including $T{=}0.05$, where every other model tested shows its pathology most strongly. This is a null result bounded by the tested budget, not evidence that the model cannot loop. An attempt to probe lower temperatures ($T\in\{0.005,0.01,0.02,0.03,0.05\}$) was abandoned before completion because of memory pressure from concurrent pilots, so no full comparison is reported for Qwen2.5-3B and no claim is made about lower temperatures, other prompt sets, or longer generations. Not every modern pretrained model exhibits this pathology at a moderate compute budget, and the reason for this one is unknown.

OLMo and Mistral both show a clear regime (peaks of 44.0\% and 20.0\%, both at $T{=}0.10$), so both proceed to full comparisons (Tables~\ref{tab:olmofull} and~\ref{tab:mistralfull}), after calibration on each model's own tokenizer, never reusing another model's frozen $(R,P_{\max})$. Neither calibration met the pre-declared 1\% tolerance, and both fall back to the smallest worst-case upper CI on the grid. OLMo freezes at $R{=}6$, $P_{\max}{=}20$ tokens (worst-case upper CI 1.38\%; validation battery 0/1391, 95\% CI [0.0\%, 0.26\%], zero false positives on every one of seven held-out corpora). Mistral freezes at $R{=}5$, $P_{\max}{=}20$ tokens (worst-case upper CI 3.15\%; validation battery 22/1400, 1.6\%, 95\% CI [1.0\%, 2.4\%], all concentrated in \emph{Dracula}'s epistolary formulas at 11.0\%, with 0.0\% on the other six corpora). The same corpus-specific false-positive concentration seen for the primary event therefore reappears on an independent tokenizer and model, which suggests it is a property of that text and not of one detector implementation.

\begin{table}[h]
\centering
\caption{OLMo-2-0425-1B (1B), $T{=}0.10$, 10 prompts $\times$ 5 seeds, 150 new tokens, frozen event ($R{=}6$, $P_{\max}{=}20$ tokens). KL against the shared $T{=}0.10$ reference, as in Table~\ref{tab:commonkl}.}
\label{tab:olmofull}
\begin{tabular}{@{}lrrrr@{}}
\toprule
Method & Loop \% & RMST (tok) & mcNLL & $\KL(q\Vert p)$ \\
\midrule
Greedy                      & 10.0\% & 144.6 & 0.999 & --- \\
Temperature 0.10 (reference) & 12.0\% & 144.5 & 0.979 & 0.000 \\
Temperature 0.05             & 14.0\% & 143.2 & 0.994 & 0.017 \\
Temperature 0.15             & 10.0\% & 146.3 & 1.037 & 0.018 \\
Repetition penalty 1.3        & 0.0\%  & 150.0 & 1.510 & 2.213 \\
No-repeat 3-gram$^{*}$       & 0.0\%  & 150.0 & 1.712 & 2.871 \\
SuffixMatch ($\alpha{=}3$)    & \textbf{0.0\%}  & 150.0 & 1.270 & \textbf{0.223} \\
FSD-style (reimpl.)          & 0.0\%  & 150.0 & 1.663 & 3.188 \\
LZ Penalty (reimpl.)        & 0.0\%  & 150.0 & 1.496 & 2.040 \\
\textsc{rr-fixed}             & 12.0\% & 146.2 & 1.083 & 0.015 \\
\textsc{rr-adaptive}          & \textbf{2.0\%}  & 149.6 & 1.070 & 0.027 \\
\textsc{rr-dual} ($\varepsilon{=}0.05$) & 0.0\% & 150.0 & 1.857 & 4.389 \\
\textsc{rr-dual} ($\varepsilon{=}0.01$) & 0.0\% & 150.0 & 1.866 & 4.590 \\
\bottomrule
\end{tabular}
\vspace{4pt}
{\footnotesize $^{*}$Hard constraint on the measured event, excluded from ranking on the same grounds as Table~\ref{tab:main}.}
\end{table}

\begin{table}[h]
\centering
\caption{Mistral-7B-v0.3 (7B), $T{=}0.10$, 10 prompts $\times$ 5 seeds, 150 new tokens, frozen event ($R{=}5$, $P_{\max}{=}20$ tokens).}
\label{tab:mistralfull}
\begin{tabular}{@{}lrrrr@{}}
\toprule
Method & Loop \% & RMST (tok) & mcNLL & $\KL(q\Vert p)$ \\
\midrule
Greedy                      & 10.0\% & 145.0 & 0.725 & --- \\
Temperature 0.10 (reference) & 4.0\%  & 148.8 & 0.727 & 0.000 \\
Temperature 0.05             & 6.0\%  & 147.2 & 0.724 & 0.013 \\
Temperature 0.15             & 10.0\% & 146.2 & 0.756 & 0.014 \\
Repetition penalty 1.3        & 0.0\%  & 150.0 & 1.976 & 8.059 \\
No-repeat 3-gram$^{*}$       & 0.0\%  & 150.0 & 1.479 & 4.500 \\
SuffixMatch ($\alpha{=}3$)    & 4.0\%  & 148.1 & 0.898 & 0.242 \\
FSD-style (reimpl.)          & \textbf{0.0\%}  & 150.0 & 1.488 & \textbf{4.252} \\
LZ Penalty (reimpl.)        & 2.0\%  & 149.4 & 0.974 & 1.004 \\
\textsc{rr-fixed}             & 4.0\%  & 148.5 & 0.791 & 0.016 \\
\textsc{rr-adaptive}          & 4.0\%  & 148.3 & 0.790 & 0.013 \\
\textsc{rr-dual} ($\varepsilon{=}0.05$) & 0.0\% & 150.0 & 1.805 & 6.739 \\
\textsc{rr-dual} ($\varepsilon{=}0.01$) & 0.0\% & 150.0 & 1.843 & 7.257 \\
\bottomrule
\end{tabular}
\vspace{4pt}
{\footnotesize $^{*}$Hard constraint on the measured event, excluded from ranking on the same grounds as Table~\ref{tab:main}.}
\end{table}

\textbf{The adaptive controller's benefit is model-dependent.} On OLMo, \textsc{rr-adaptive} cuts the loop rate from the 12.0\% shared by the baseline and \textsc{rr-fixed} to 2.0\% at negligible KL cost (0.027 bits). This is not an artifact of the calibrated threshold: across a full $(R,P_{\max})$ sensitivity sweep (5 values of $R\in\{4,\ldots,8\}\times5$ values of $P_{\max}\in\{10,\ldots,40\}$, 25 grid points, using the raw generated token ids saved for this purpose), the ordering $\textsc{rr-adaptive}\le\textsc{rr-fixed}\le$ baseline holds at every grid point, with 0 reversals out of 25. On Mistral, \textsc{rr-adaptive} and \textsc{rr-fixed} are indistinguishable from the $T{=}0.10$ baseline at the frozen threshold (4.0\% for all three), and the same sweep shows no consistent direction: at looser thresholds ($R{=}3,4$) both sit below the baseline, at the calibrated and stricter thresholds ($R{=}5$--$7$, tight $P_{\max}$) at or slightly above it (e.g.\ $R{=}5$, $P_{\max}{=}10$: baseline 0.0\%, both variants 4.0\%). The same mechanism and hyperparameters help substantially on one 2024-generation checkpoint and do nothing distinguishable from noise on another; I have no explanation for this.

\textbf{The cheapest route to zero measured loops differs by model.} On OLMo, SuffixMatch is the cheapest of the zero-loop methods (0.223 bits), whereas on the primary model LZ Penalty is (Table~\ref{tab:commonkl}). GPT-2 has no comparable KL ranking (loop rate, RMST, and mcNLL only), and SuffixMatch does not reach zero loops there (4.0\%). On Mistral, neither SuffixMatch (4.0\%) nor LZ Penalty (2.0\%) reaches zero loops; FSD-style is the only external baseline that does, at 4.252 bits. \textsc{rr-dual} remains the most expensive route to zero loops on both models (4.39--4.59 bits on OLMo, 6.74--7.26 bits on Mistral). The detector-threshold sweep covered a seven-configuration subset (the recurrence-risk family, SuffixMatch, and FSD-style); on OLMo, SuffixMatch and both dual configurations stay at exactly 0.0\% at every one of the 25 grid points, and on Mistral SuffixMatch's loop rate ranges 4.0--20.0\% across the 25 cells, consistent with the 4.0\% at the calibrated point. The sweep for FSD-style used an earlier implementation and is not reported, and LZ Penalty was not included.

\textbf{Generation length does not hide a different picture at 150 tokens.} Re-censoring each sample's saved onset at shorter horizons (30, 60, 90, 120 tokens) shows OLMo's \textsc{rr-adaptive} advantage opening in the back half of generation (0.0\% for both \textsc{rr-adaptive} and baseline through 90 tokens; the baseline reaches 10--12\% by 150 tokens against 0--2\% for \textsc{rr-adaptive}). The Mistral null finding holds at every horizon tested.

\textsc{rr-dual} reaches exactly 100.0\% per-step feasibility on both OLMo and Mistral (both $\varepsilon$, zero structurally infeasible and zero near-boundary steps), so the guarantee has held without exception on every model tested: the primary character model and its second checkpoint, the non-fiction model, GPT-2, OLMo, and Mistral. Together with the GPT-2 character-level $(R,P_{\max})$ sensitivity check (25 grid points, all zero-loop findings unchanged), this shows that a mechanism's effectiveness, unlike the dual mechanism's per-step feasibility, is an empirical property of a specific model and cannot be assumed to transfer from one modern checkpoint to the next.

%==============================================================================
\section{Discussion}
\label{ssec:discussion}
%==============================================================================

\subsection{What is supported}
\textbf{The local surrogate problem's exact solution behaves as the theorem says.} On the validated event, \textsc{rr-dual} eliminates persistent loops on the primary model and on the independently trained non-fiction model at a corpus-consistent cost (11.2 and 12.9 bits at $\varepsilon{=}0.05$), and its per-step feasibility is exactly 100.0\% on every model tested. This is the strongest generalization evidence in the paper for the local-surrogate claim, because it follows from the theory (exact per-step recalibration) and not from corpus-specific tuning. \textsc{rr-fixed} reduces the loop rate substantially on the primary model at very low cost, but that effectiveness does not transfer to the second corpus without re-tuning: exactness and effectiveness are different axes, and only the surrogate problem's exactness generalizes cleanly.

\textbf{The gap to path-level event control is measured exactly.} On the synthetic problem, hazard-aware local surrogate control has the lower loop rate than local-only control at all 5 grid points where a matched-KL comparison is possible; at the tightest budget local's structural KL ceiling ($\approx0.613$ bits/step) makes the comparison inapplicable. Both are local surrogate mechanisms, and both fall well short of the path-level oracle at every budget (1.6--5.1$\times$ its KL, Section~\ref{ssec:oracle}). Hazard-aware control's edge over plain local control is an edge between two surrogates and not a step toward closing the surrogate-to-path-level gap, which the oracle's closed form shows is a matter of how the hazard signal is used (exponentially, not affinely) and not only of how coarse the signal is.

\subsection{What is not supported}
The following concerns which local surrogate mechanism to prefer and whether a mechanism's effectiveness generalizes, not path-level control, which no mechanism here attempts on a real model.
\begin{itemize}[leftmargin=*]
  \item Recurrence-risk decoding is not shown to improve survival over baselines at ordinary sampling temperatures, where the primary model does not exhibit the pathology, and \textsc{rr-fixed} is not shown to beat repetition penalty on survival on the primary model ($p{=}0.246$).
  \item The hazard-aware mechanism is not shown to transfer to a real model (Section~\ref{sec:hazard-transfer}); the null result is specific to the rollout estimator tested.
  \item The effectiveness of the heuristic configurations is model-specific. LZ Penalty is 0\% on the primary model, 2.0\% on the non-fiction corpus, and 8.0\% on GPT-2; FSD-style is 1.3\% and 16.7\% on the first two and 0.0\% on GPT-2. The adaptive controller's benefit on OLMo (12.0\%$\to$2.0\%, robust across a 25-point detector-threshold sweep) is indistinguishable from noise on Mistral, and the cheapest zero-loop baseline differs by model (LZ Penalty on the primary model, SuffixMatch on OLMo, FSD-style on Mistral). Qwen2.5-3B shows no instance of the pathology at the budget tested.
  \item The human evaluation ($n{=}9$) detects no reliable preference ranking among the zero-loop group, and is not evidence that the decoders are equally readable.
\end{itemize}
The dual guarantee is a one-step bound on the \emph{expected} recurrence risk, given that the numerical projection succeeds; it does not bound the sequence-level probability of ever looping, and every 0/150 figure in this paper is an observed count, not a proven zero probability.

\subsection{Limitations}
\begin{itemize}[leftmargin=*]
  \item \textbf{Detector.} The pooled false-positive rate of $\sim$1.3\% conceals per-corpus variance: \emph{Dracula} alone fires at 26/300 = 8.7\% (Table~\ref{tab:validation}), about 7$\times$ the pooled figure, concentrated in epistolary formulas that a syntactic event cannot distinguish from degeneration. The non-fiction corpus's $\sim$0.8\% was not checked for single-corpus concentration. The event is syntactic and does not capture semantic repetition, and mcNLL is an in-distribution proxy, not a validated quality measure.
  \item \textbf{Solver.} Residual dual infeasibility is exactly $0.0\%$ on both character-model checkpoints. It could still arise on another model or risk configuration, which the closed-form check (Section~\ref{ssec:dual-fix}) would detect.
  \item \textbf{Baselines.} FSD-style implements the published equations \citep{yang2023frustratingly} and is not a certified reproduction, since no reference code is available. LZ Penalty is ported from the authors' reference and verified against it. \textsc{SuffixMatchDecoder} is a homemade heuristic, not the published Look-back algorithm.
  \item \textbf{Generalization.} The non-fiction and GPT-2 comparisons did not repeat the temperature sweep, the repetition-penalty sweep, or fixed-$\lambda$ re-tuning that might restore \textsc{rr-fixed} and \textsc{rr-adaptive} there. GPT-2's risk order ($n_{\min}{=}2,n_{\max}{=}4$) and LZ Penalty buffer/window (32/512) were set once from published subword defaults and not swept. The common-distortion ranking (Table~\ref{tab:commonkl}) is a single-checkpoint, single-temperature snapshot.
  \item \textbf{Human evaluation.} Participants were recruited informally, and all sessions used the primary character model at $T{=}0.10$.
  \item \textbf{Modern models.} The comparison is compact (5 seeds $\times$ 10 prompts $\times$ 150 tokens per configuration) and covers only $T{=}0.10$, so confidence intervals are wide. Neither calibration met the pre-declared 1\% tolerance (OLMo 1.38\%, Mistral 3.15\% worst-case upper CI). Threshold sensitivity was checked for a seven-configuration subset that excludes LZ Penalty, and FSD-style's has not been recomputed with the current implementation; generation-length sensitivity was checked for all fourteen configurations.
\end{itemize}

%==============================================================================
\section{Conclusion}
%==============================================================================

Local surrogate control and path-level event control are different problems. Recurrence-risk decoding, dual calibration included, solves only the first exactly, and the distance to the second is measured here exactly. On a finite-state problem where the path-level optimum is computable, even a one-step projection given the true multi-step hazard spends on average 2.2$\times$ the KL of the optimal path policy, and a plain local surrogate 3.1$\times$. The oracle's one-step marginal reweights the hazard affinely, whereas hazard-aware control reweights it exponentially, so the gap is not closed by a better risk signal.

On real models, recurrence-risk decoding removes a validated failure mode at an exactly quantified cost. Its dual mode guarantees a per-step bound on expected recurrence risk whenever the numerical projection succeeds; that success rate is exactly 100.0\% on every model tested, and zero persistent loops was observed on the primary and non-fiction models, GPT-2, OLMo, and Mistral as an empirical finding, not a consequence of the theorem. The cheaper heuristic variants do not carry their effectiveness across models: \textsc{fixed} ranges from 2.7\% to 51.3\% loops across corpora, and \textsc{adaptive} helps on OLMo and not on Mistral. Against faithful reimplementations of FSD and the LZ Penalty, LZ Penalty is the cheapest external baseline that reaches zero loops on the primary model (0.186 bits) but fails to reach zero on the non-fiction corpus and on GPT-2, while FSD-style fails to reach zero on the primary model. \textsc{rr-dual} is the most expensive route to zero loops on every model tested, yet a blinded human evaluation with nine participants finds no reliable readability difference among the zero-loop decoders. The hazard-aware extension is validated where the path-level answer is exactly known and yields no detectable signal on a real model under the rollout estimator tested.

Closing the gap to path-level event control remains open. Neither the guaranteed nor the heuristic one-step mechanisms tested here address it, and the oracle's closed-form marginal suggests where a tractable approximation might start.

\bibliographystyle{plainnat}
\bibliography{references}

\appendix

\section{Implementation Defects and Corrections}
\label{app:defects}
This appendix records defects found during development, their effect on reported numbers, and their corrections. Every table in the paper reflects the corrected code.

\subsection*{Loop event and reproducibility}
\begin{itemize}[leftmargin=*]
  \item \textbf{Window-sampling seed.} The calibration and validation script seeded its window sampling with Python's built-in \texttt{hash()} on corpus names, which is randomized per process, so the sampled windows differed between invocations. It now uses a seed derived from \texttt{hashlib.md5}, verified identical across processes, and saves the sampled window start indices in its report. Re-running reproduces the frozen setting ($R{=}5$, $P_{\max}{=}60$) and closely similar false-positive counts.
  \item \textbf{Period-1 threshold.} The event was first described as ``$R{=}5$'' without qualification; period 1 uses a decoupled threshold of 8 (Section~\ref{sec:event}).
  \item \textbf{Comparison keys.} \texttt{compute\_\allowbreak{}significance.py} used the pre-rename key \texttt{lt\_lookback\_a3.0}, which silently dropped that comparison from every run after the rename to \texttt{lt\_suffixmatch\_a3.0}. The class \texttt{LookBackDecoder} was renamed \textsc{SuffixMatchDecoder} because the name implied it was the published Look-back algorithm. Greedy decoding was initially evaluated on 5 of the 15 test prompts and is now evaluated on all 15.
\end{itemize}

\subsection*{Model-consistency NLL}
The first implementation scored non-overlapping chunks of length \texttt{block\_\allowbreak{}size}. On a toy model, about half of all scored positions received less context than the architecture allows (in one case 9 of 16 possible context tokens). A first fix (overlapping windows with stride $=$ block\_size$/2$, keeping each window's back half) was still insufficient, since only the last position of each window reached full context. The final method scores every target position independently, batching each position's exact context window; on the toy model all 40 scored positions receive exactly $\min(t,\text{block\_size})$ tokens (0 mismatches). An intermediate version left-padded windows with a real vocabulary token id, which unmasked causal attention legitimately attended to, corrupting the scored logits (0.26 maximum logit difference between isolated and padded evaluation of the same context). Right-padding is used instead, since causal masking cannot let it leak into any real position, and the batched computation matches an isolated forward pass to 6 decimals.

\subsection*{Dual solver}
\begin{enumerate}[leftmargin=*]
  \item \textbf{Float32 underflow.} The solver computed $q_\lambda$'s logits as $\log(p+10^{-40})-\lambda\risk$ from a materialized softmax. At $T{=}0.10$, 276 of 2000 tested vocabulary entries underflowed to exactly $0.0$ in one constructed check, with true log-probabilities spanning 41 nats ($-104$ to $-145$) collapsed to a single floor value ($\log 10^{-40}\approx-92.1$). Since dual calibration reweights exactly this tail, this changed $\lambda$, $q$, and KL. The solver now works in log-space throughout.
  \item \textbf{Infeasibility conflated with the cap.} True infeasibility ($\varepsilon<\min_v\risk(v)$) was learned only by bisecting to a fixed cap and failing. It is now decided in closed form before any search, and the search bracket expands geometrically (Section~\ref{ssec:dual-fix}); the earlier caps ($\lambda_{\max}{=}50$, later 1000) are removed. An earlier description of the search as having ``no ceiling'' was incorrect: doubling stops after 60 iterations.
  \item \textbf{Discarded diagnostics.} An earlier run hid a solver that was silently infeasible on $\sim$26--28\% of steps.
  \item \textbf{Float32/float64 comparison at the boundary.} At exactly $\varepsilon=\min_v\risk(v)$ the check between a float32 \texttt{risk} tensor and a float64 \texttt{eps} could round the wrong way and misreport a feasible point as infeasible; it now uses a tolerance and a \texttt{near\_\allowbreak{}boundary} flag.
  \item \textbf{NaN diagnostics.} A NaN entered the KL and entropy diagnostics whenever top-$p$ filtering ran after the projection with $q(v){=}0$ exactly.
\end{enumerate}
\textbf{Effect.} Dual KL on the primary model was 2.256--2.328 bits in the original run, 10.3--11.6 bits after the cap was fixed but the float32 arithmetic was not, and 11.2--12.7 bits after the numerical fix, about five times the original figure. Per-step feasibility was 72.3--74.5\% originally, 98.75--98.87\% after the cap fix, and is exactly 100.0\% at both $\varepsilon$ after the numerical fix. Loop rate and RMST do not depend on this arithmetic and were unaffected.

\subsection*{Theory}
The theorem originally omitted the $\varepsilon=\min_v\risk(v)$ boundary, where no finite dual variable exists. The claim that the method ``becomes standard hard no-repeat'' as $\lambda\to\infty$ was wrong twice: it fails when $n_{\min}<n_{\max}$, and even for a single order it ignores that hard blocking crashes when no safe continuation exists (Section~\ref{sec:theory}).

\subsection*{Synthetic experiments}
\begin{itemize}[leftmargin=*]
  \item \textbf{Unreproducible headline.} An earlier version reported ``17.9\% vs.\ 53.2\% loop rate at matched KL$\approx0.24$'', which could not be reproduced from its saved artifact: the saved check targeted a different operating point than the one cited. The matched-distortion comparison is now computed and saved for every grid point.
  \item \textbf{Mismatched $\varepsilon$.} Local and hazard-aware control were called with a one-step $\varepsilon$ read directly off the grid used for the oracle's sequence-level target, so the three mechanisms were compared at different operating points. All three are now matched at the same sequence-level $\varepsilon$ (Section~\ref{ssec:oracle}); this changed the number of comparable grid points and the reported margins (for example, at $\varepsilon{=}0.20$ the hazard loop rate is now exactly 20.0\%, replacing a simulated 17.9\%).
  \item \textbf{Ignored infeasibility flag.} The one-step solver's structural-infeasibility flag was computed but never checked; it is now tracked and reported.
  \item \textbf{Non-absorbing KL.} The local and hazard simulator kept charging KL for the full nominal horizon after a trajectory had looped, inflating KL/step. KL is now absorbing at the loop event, matching the oracle. This changed the matched-local $\varepsilon$ and loop-rate columns of Table~\ref{tab:synthetic} (for example at $\varepsilon{=}0.20$, local loop rate from 51.6\% to 44.2\%); Table~\ref{tab:oracle} was unaffected, since it already used the exact absorbing forward propagation. Margins changed from 7.8--37.6 to 7.6--31.1 percentage points.
  \item \textbf{Base of the exponent.} The path-level derivation mixed $2^{-\lambda}$ and $e^{-\lambda}$; it now uses base 2 throughout, and the closed-form identity holds to $<10^{-14}$.
\end{itemize}

\subsection*{Baseline reimplementations}
The FSD-style and LZ Penalty baselines were reimplemented to match their published specifications. The earlier FSD-style combined the model and anti-LM in logit space with a longest-available-order heuristic and scored the full vocabulary; the earlier LZ Penalty scored single tokens, used a fallback formula that differs from the paper's, and had the wrong sign convention. Effects of the corrections on reported figures: primary-model FSD-style loop rate 0.0\%$\to$1.3\%; LZ Penalty is now the cheapest external zero-loop baseline on the primary model (0.186 bits); non-fiction FSD-style 0.0\%$\to$16.7\% and LZ Penalty 0.0\%$\to$2.0\%; GPT-2 FSD-style 1.0\%$\to$0.0\% and LZ Penalty 25.0\%$\to$8.0\%; Mistral FSD-style KL 1.691$\to$4.252 bits and LZ Penalty loop rate 6.0\%$\to$2.0\%; OLMo FSD-style KL 1.510$\to$3.188 bits and LZ Penalty KL 2.124$\to$2.040 bits; spelling-error rates FSD-style 4.28\%$\to$2.59\%, LZ Penalty 1.58\%$\to$0.75\%. The paired-comparison rows involving FSD-style were recomputed (\texttt{compute\_\allowbreak{}significance.py} now includes them). The human-evaluation passages were regenerated from the corrected decoders and the evaluation collected anew; sessions collected on the earlier passages are superseded.

\subsection*{Detector-threshold sensitivity}
The $(R,P_{\max})$ sweep for FSD-style on Mistral was run with the earlier implementation (12.0\% loops at $R{=}3$, 0.0\% for $R\ge4$) and has not been recomputed; it is not used for any claim in the main text.

\section{Order-of-Magnitude Estimate for the Multiplier Cap}
\label{app:lambda-estimate}
For a two-outcome simplification (one risky token with $\risk{=}1$, one safe token with $\risk{=}0$, raw pre-temperature probability ratio $\rho$ risky:safe), the $\lambda$ needed to bring $q(\text{risky})\le\varepsilon$ at temperature $T$ is $\lambda\ge\ln\frac{1-\varepsilon}{\varepsilon}+\frac{1}{T}\ln\rho$. At $T{=}0.10$, $\rho{=}100$ needs $\lambda\approx50$ and $\rho{=}1000$ needs $\lambda\approx73$; a cap of 1000 already covers any $\rho$ up to $\ln\rho\lesssim97$. This estimate is for an interior point, away from the $\varepsilon=\min_v\risk(v)$ boundary; at that boundary the required $\lambda$ depends on how small a float64 gap the solver must resolve against a float32-rounded risk floor, not on a probability ratio, and is handled by the \texttt{near\_\allowbreak{}boundary} tolerance. The 60-doubling limit ($\lambda\le2^{60}\approx1.15\times10^{18}$) exceeds the $\rho{=}1000$ estimate by roughly 16 orders of magnitude and is not observed to bind outside that boundary case.

\section{Reproducibility Bundle}
\label{sec:reproducibility}
\begin{itemize}[leftmargin=*]
  \item \textbf{Decoders and evaluation.} \texttt{src/decoding.py} (all decoders, the log-space dual solver with its feasibility and cap diagnostics, and opt-in per-step diagnostics via \texttt{per\_\allowbreak{}step\_\allowbreak{}log}); \texttt{sampling\_\allowbreak{}eval.py} (with \texttt{--data\_\allowbreak{}dir}, \texttt{--ckpt}, and \texttt{--save\_\allowbreak{}per\_\allowbreak{}step\_\allowbreak{}diagnostics}); \texttt{merge\_\allowbreak{}sampling\_\allowbreak{}eval\_\allowbreak{}runs.py} (recombines partial \texttt{--only} reruns); \texttt{compute\_\allowbreak{}significance.py} (source of Table~\ref{tab:paired}); \texttt{common\_\allowbreak{}kl\_\allowbreak{}comparison.py} (Table~\ref{tab:commonkl}); \texttt{temperature\_\allowbreak{}sweep.py}.
  \item \textbf{Loop events.} \texttt{validate\_\allowbreak{}loop\_\allowbreak{}event\_\allowbreak{}v2.py} (deterministic seed; \texttt{--confirm\_\allowbreak{}only} mode for the second corpus); \texttt{validate\_\allowbreak{}loop\_\allowbreak{}event\_\allowbreak{}subword.py}; \texttt{loop\_\allowbreak{}event\_\allowbreak{}subword\_\allowbreak{}report\_\allowbreak{}\{olmo,mistral\}/}.
  \item \textbf{Pretrained models.} \texttt{pretrained\_\allowbreak{}subword\_\allowbreak{}pilot.py}, \texttt{gpt2\_\allowbreak{}full\_\allowbreak{}comparison.py}, and \texttt{modern\_\allowbreak{}llm\_\allowbreak{}comparison.py} (pilot, calibrate, compare on OLMo-2-0425-1B, Qwen2.5-3B, and Mistral-7B-v0.3; KV-cached generation; \texttt{--only} for a configuration subset; raw \texttt{gen\_\allowbreak{}ids} saved per sample to support threshold sensitivity sweeps without regeneration); \texttt{sensitivity\_\allowbreak{}pareto\_\allowbreak{}analysis.py} (generation-length and detector-threshold sensitivity from saved artifacts).
  \item \textbf{Synthetic and hazard experiments.} \texttt{synthetic\_\allowbreak{}fsm\_\allowbreak{}hazard.py} (whole-grid matched-distortion computation and the exact sequence-level oracle); \texttt{build\_\allowbreak{}hazard\_\allowbreak{}dataset.py}; \texttt{test\_\allowbreak{}multistep\_\allowbreak{}hazard.py}.
  \item \textbf{Human evaluation.} \texttt{human\_\allowbreak{}eval\_\allowbreak{}tool.html} and \texttt{aggregate\_\allowbreak{}human\_\allowbreak{}eval.py}; \texttt{human\_\allowbreak{}eval\_\allowbreak{}results\_\allowbreak{}v2/} holds the raw exports of the 9 participants analysed in Section~\ref{sec:human-eval}, and \texttt{human\_\allowbreak{}eval\_\allowbreak{}results/} the superseded exports on an earlier passage set.
  \item \textbf{Result files.} \texttt{all\_\allowbreak{}results.json}, \texttt{metrics\_\allowbreak{}*.csv}, \texttt{synthetic\_\allowbreak{}fsm\_\allowbreak{}results.json}, \texttt{significance\_\allowbreak{}table7.\{json,csv\}}, \texttt{samples\_\allowbreak{}*.txt}, \texttt{loop\_\allowbreak{}event\_\allowbreak{}confirm\_\allowbreak{}report.json}, \texttt{loop\_\allowbreak{}event\_\allowbreak{}subword\_\allowbreak{}report.json}, \texttt{pretrained\_\allowbreak{}pilot\_\allowbreak{}results.json}, and \texttt{per\_\allowbreak{}step\_\allowbreak{}diagnostics/*.jsonl} (dual-mode configurations, primary model).
\end{itemize}

\section{Prompt Sets}
\textbf{\textsc{Dev} (5, tuning only):} \emph{``CHAPTER 1''}; \emph{``The night was''}; \emph{``She had never''}; \emph{``It was the best of''}; \emph{``Mr.\ Darcy had never''}.
\textbf{\textsc{Test} (15, untouched during tuning, used identically for both models):} \emph{``The road turned sharply where''}; \emph{``He folded the letter and''}; \emph{``Rain struck the window as''}; \emph{```You must understand,' she began,''}; \emph{``At the market that morning''}; \emph{``Three years afterwards, when''}; \emph{``The room smelled of''}; \emph{``The doctor set down his bag and''}; \emph{``Beyond the river lay''}; \emph{```And what,' he asked, `became of''}; \emph{``That winter the family''}; \emph{``The ship had been at sea for''}; \emph{``In the garden behind the house''}; \emph{``His debts had grown until''}; \emph{``By dawn the streets were''}.

\end{document}
